\lecture{17}
\title{2D Fourier Transform I}

\begin{document}

\subsection{Lecture: 2D Fourier Series}

For constructing 1D Fourier series, the main idea was that
for a given period $T$ and fundamental frequency $\omega_0 = \frac{2\pi}{T}$,
the complex exponentials form an orthogonal basis.
\[ \phi_{k}(t) := e^{jk \omega_0 t}
\implies
\left \langle \phi_{k_1}, \phi_{k_2} \right \rangle
= \int_{T} \left[e^{-jk_1 \omega_0 t}\right] \cdot \left[e^{jk_2 \omega_0 t}\right]
\, \mathrm{d}t
= \begin{cases}
T & \text{ if } k_1 = k_2 \\
0 & \text{ otherwise.}
\end{cases} \]
This generalizes readily to 2D Fourier series.
We have a signal $f: \mathbb{R} \times \mathbb{R} \to \mathbb{C}$
with periods $T_x$ and $T_y$ along the $x$ and $y$ directions,
with directional fundamental frequencies
$\omega_{0, x} = \frac{2\pi}{T_x}$ and $\omega_{0, y} = \frac{2\pi}{T_y}$.
Then our orthogonal basis is:
\begin{align*}
    \phi_{k_x, k_y}(x, y) := e^{jk_x \omega_{0, x} x} \cdot e^{jk_y \omega_{0, y} y}
    ~ \implies ~~~~~~
    & \left \langle \phi_{k_{x, 1}, k_{y, 1}}, \phi_{k_{x, 2}, k_{y, 2}} \right \rangle \\
    = \ & \int_{T_x} \int_{T_y} \left[e^{-jk_{x, 1} \omega_{0, x} x} \cdot e^{-jk_{y, 1} \omega_{0, y} y}\right]
    \cdot \left[e^{jk_{x, 2} \omega_{0, x} x} \cdot e^{jk_{y, 2} \omega_{0, y} y}\right]
    \, \mathrm{d}y \, \mathrm{d}x \\
    = \ & \begin{cases}
    T_x T_y & \text{ if } (k_{x, 1}, k_{y, 1}) = (k_{x, 2}, k_{y, 2}). \\
    0 & \text{ otherwise.}
    \end{cases}
\end{align*}

By corollary, 2D Fourier series work like so:

\begin{theorem}{2D CTFS}
    Given a periodic signal $f: \mathbb{R} \times \mathbb{R} \to \mathbb{C}$,
    its Fourier series works like so.
    \[ F[k_x, k_y] := \frac{1}{T_xT_y} \int_{T_x} \int_{T_y}
    f(x, y) \left[e^{-jk_x \omega_{0, x} x} \cdot e^{-jk_y \omega_{0, y} y}\right]
    \, \mathrm{d}y \, \mathrm{d}x
    ~~~ \parallel ~~~
    f(x, y) = \sum_{k_x = -\infty}^{\infty} \sum_{k_y = -\infty}^{\infty}
    F[k_x, k_y] \left[e^{jk_x \omega_{0, x} x} \cdot e^{jk_y \omega_{0, y} y}\right]. \]
\end{theorem}

The 2D DTFS works exactly the same.

\begin{theorem}{2D DTFS}
    Given a periodic signal $f: \mathbb{Z} \times \mathbb{Z} \to \mathbb{C}$,
    its Fourier series works like so.
    \[ F[k_x, k_y] := \frac{1}{N_xN_y} \sum_{n_x = 0}^{N_x - 1} \sum_{n_y = 0}^{N_y - 1}
    f[n_x, n_y] \left[e^{-jk_x \Omega_{0, x} n_x} \cdot e^{-jk_y \Omega_{0, y} n_y}\right]
    ~~~ \parallel ~~~
    f[n_x, n_y] = \sum_{k_x = 0}^{N_x - 1} \sum_{k_y = 0}^{N_y - 1}
    F[k_x, k_y] \left[e^{jk_x \Omega_{0, x} n_x} \cdot e^{jk_y \Omega_{0, y} n_y}\right]. \]
\end{theorem}

\subsection{Lecture: 2D Fourier Transform}

Just as before, we care more about the CTFT and DTFT
than we do the CTFS and DTFS.
Their derivations work just as before.

\begin{theorem}{2D CTFT}
    Given a signal $f: \mathbb{R} \times \mathbb{R} \to \mathbb{C}$,
    its Fourier transform is:
    \[ F(\omega_x, \omega_y) :=
    \int_{-\infty}^{\infty} \int_{-\infty}^{\infty}
    f(x, y) e^{-j(\omega_x x + \omega_y y)}
    \, \mathrm{d}x \, \mathrm{d}y
    ~~~ \parallel ~~~
    f(x, y) = \frac{1}{4\pi^2}
    \int_{-\infty}^{\infty} \int_{-\infty}^{\infty}
    F(\omega_x, \omega_y) e^{j(\omega_x x + \omega_y y)}
    \, \mathrm{d}\omega_x \, \mathrm{d}\omega_y. \]
\end{theorem}

\begin{proof}
    As in \href{6.html#lecture-the-ct-fourier-transform}{Lecture 6},
    take $f_{T_x, T_y}(x, y)$ to be the extension of $f(x, y)$
    to have period $(T_x, T_y)$.  The CTFS of $f_{T_x, T_y}(x, y)$ is:
    \[ F_{T_x, T_y}[k_x, k_y] = \frac{1}{T_xT_y}
    \int_{T_x} \int_{T_y}
    f(x, y) \left[e^{-j\textcolor{purple}{k_x \omega_{0, x}} x}
    \cdot e^{-j\textcolor{purple}{k_y \omega_{0, y}} y}\right]
    \, \mathrm{d}y \, \mathrm{d}x \]
    Now define the CTFT to be $F(\omega_x, \omega_y) :=
    \lim_{(T_x, T_y) \to (\infty, \infty)} T_xT_y \cdot F_{T_x, T_y}[k_x, k_y]$,
    where \textcolor{purple}{$\omega_x = k_x \omega_{0, x}$} and
    \textcolor{purple}{$\omega_y = k_y \omega_{0, y}$} are the reparametrization.
    First, check this matches the theorem statement.
    \[ F(\omega_x, \omega_y)
    = \lim_{(T_x, T_y) \to (\infty, \infty)}
    T_xT_y \cdot \left[
        \frac{1}{T_xT_y}
        \int_{T_x} \int_{T_y}
        f(x, y) \left[e^{-j\textcolor{purple}{k_x \omega_{0, x}} x}
        \cdot e^{-j\textcolor{purple}{k_y \omega_{0, y}} y}\right]
        \, \mathrm{d}y \, \mathrm{d}x
    \right]
    = \int_{-\infty}^{\infty} \int_{-\infty}^{\infty}
    f(x, y) e^{-j(\textcolor{purple}{\omega_x} x + \textcolor{purple}{\omega_y} y)}
    \, \mathrm{d}x \, \mathrm{d}y. \]
    Next, see what happens to the CTFS synthesis equation
    at the limit.
    \begin{align*}
    f(x, y) & = \lim_{(T_x, T_y) \to (\infty, \infty)}
    \sum_{k_x = -\infty}^{\infty} \sum_{k_y = -\infty}^{\infty}
    F_{T_x, T_y}[k_x, k_y]
    \left[e^{j\textcolor{purple}{k_x \omega_{0, x}} x} \cdot
    e^{j\textcolor{purple}{k_y \omega_{0, y}} y}\right] \\
    & = \lim_{(T_x, T_y) \to (\infty, \infty)}
    \sum_{k_x = -\infty}^{\infty} \sum_{k_y = -\infty}^{\infty}
    \textcolor{red}{\left[T_xT_y \cdot F_{T_x, T_y}[k_x, k_y]\right]}
    \left[e^{j\textcolor{purple}{\omega_{k_x}} x} \cdot
    e^{j\textcolor{purple}{\omega_{k_y}} y}\right]
    \left[\frac{1}{4\pi^2} \cdot
    \textcolor{red}{\frac{2\pi}{T_x}} \cdot
    \textcolor{red}{\frac{2\pi}{T_y}}\right] \\
    & = \lim_{(T_x, T_y) \to (\infty, \infty)}
    \sum_{k_x = -\infty}^{\infty} \sum_{k_y = -\infty}^{\infty}
    \textcolor{red}{F(\omega_{k_x}, \omega_{k_y})}
    \left[e^{j\textcolor{purple}{\omega_{k_x}} x} \cdot
    e^{j\textcolor{purple}{\omega_{k_y}} y}\right]
    \left[\frac{1}{4\pi^2} \cdot
    \textcolor{red}{\Delta \omega_{k_x}} \cdot
    \textcolor{red}{\Delta \omega_{k_y}}\right] \\
    & = \frac{1}{4\pi^2}
    \int_{-\infty}^{\infty} \int_{-\infty}^{\infty}
    F(\omega_x, \omega_y) e^{j(\omega_x x + \omega_y y)}
    \, \mathrm{d}\omega_x \, \mathrm{d}\omega_y.
    \end{align*}
    To clarify: in the above,
    $\{\omega_{k_x}\}_{k_x = -\infty}^{\infty}$
    denotes an arithmetic progression with common difference
    $\Delta \omega_{k_x} = \omega_{0, x} = \frac{2\pi}{T_x}$.
    ~ $\blacksquare$
\end{proof}

Analogously, you can take the 2D DTFT as well.

\begin{theorem}{2D DTFT}
    Given a signal $f: \mathbb{Z} \times \mathbb{Z} \to \mathbb{C}$,
    its Fourier transform is:
    \[ F(\Omega_x, \Omega_y) :=
    \sum_{n_x = -\infty}^{\infty} \sum_{n_y = -\infty}^{\infty}
    f[n_x, n_y] e^{-j(\Omega_x n_x + \Omega_y n_y)}
    ~~~ \parallel ~~~
    f[n_x, n_y] = \frac{1}{4\pi^2}
    \int_{0}^{2\pi} \int_{0}^{2\pi}
    F(\Omega_x, \Omega_y) e^{j(\Omega_x n_x + \Omega_y n_y)}
    \, \mathrm{d}\Omega_x \, \mathrm{d}\Omega_y. \]
\end{theorem}

And of course, there's a 2D DFT as well, and it works just as before.

\begin{theorem}{2D DFT}
    Given a DT signal $f[n_x, n_y]$ of size $N_x \times N_y$,
    its discrete Fourier transform is:
    \[ F[k_x, k_y] := \frac{1}{N_xN_y} \sum_{n_x = 0}^{N_x - 1} \sum_{n_y = 0}^{N_y - 1}
    f[n_x, n_y] \left[e^{-jk_x \Omega_{0, x} n_x} \cdot e^{-jk_y \Omega_{0, y} n_y}\right]
    ~~~ \parallel ~~~
    f[n_x, n_y] = \sum_{k_x = 0}^{N_x - 1} \sum_{k_y = 0}^{N_y - 1}
    F[k_x, k_y] \left[e^{jk_x \Omega_{0, x} n_x} \cdot e^{jk_y \Omega_{0, y} n_y}\right]. \]
\end{theorem}

\begin{remark}
    The 2D DFT is equivalent to the evaluation of the polynomial
    $p(z_x, z_y) := \sum_{n_x = 0}^{N_x - 1} \sum_{n_y = 0}^{N_y - 1}
    f[n_x, n_y] z_x^{n_x} z_y^{n_y}$ as $z_x$ and $z_y$
    span the $N_x^{\text{th}}$ and $N_y^{\text{th}}$ roots of unity, respectively.
    And naturally, 2D circular convolution is just 
    multivariate polynomial multiplication in the ring
    $\mathbb{C}[z_x, z_y]/((z_x^{N_x} - 1)(z_y^{N_y} - 1))$, and so on\dots
\end{remark}

\subsection{Recitation and Lecture: 2D DFT Algebraic Properties}

\begin{theorem}{2D Unit-Sample DFT}
    The DFT of $f[n_x, n_y] := \delta[n_x] \delta[n_y]$ is
    the constant function $F[k_x, k_y] = \frac{1}{N_xN_y}$.
\end{theorem}

\begin{proof}
    Compare with the 1D case discussed in
    \href{7.html#lecture-and-recitation-properties-of-the-dt-fourier-transform}{Lecture 7}.
    Proof by plain computation. ~ $\blacksquare$
\end{proof}

Here's a useful trick for handling higher-dimension DFTs.

\begin{theorem}{2D DFT Independence}
    To compute the size-$(N_x, N_y)$ DFT of $f[n_x, n_y]$,
    it is equivalent to do the following.
    \begin{itemize}
        \item First, replace every row
        with the size-$N_x$ DFT of that row.
        \item Afterwards, replace every column
        with the size-$N_y$ DFT of that column.
    \end{itemize}
\end{theorem}

\begin{proof}
    Move the terms around between the double-summation.
    \[ F[k_x, k_y] =
    \overbrace{\frac{1}{N_y} \sum_{n_y = 0}^{N_y - 1} \underbrace{\left(
        \frac{1}{N_x} \sum_{n_x = 0}^{N_x - 1}
        f[n_x, n_y] e^{-jk_x\Omega_{0, x}n_x}
    \right)}_{
        \text{row-wise DFT}
    } e^{-jk_y\Omega_{0, y} n_y}}^{\text{column-wise DFT}}. \]
    Alternatively, adopt the polynomial interpretation:
    a multivariate polynomial in $(z_x, z_y)$
    is just a single-variable polynomial in $z_y$,
    where the coefficients are themselves polynomials in $z_x$.
    ~ $\blacksquare$
\end{proof}

Here's a visual of how the trick above
can be used to calculate the DFT of $f[n_x, n_y] = \delta[n_x]\delta[n_y]$.

\begin{center}
    \frame{\includegraphics[width=0.4\textwidth]{lecture-17/unit-sample-DFT}}
\end{center}

You can check for yourself the following facts, too:
\begin{itemize}
    \item The DFT of the constant signal $f[n_x, n_y] := 1$
    is $F[k_x, k_y] = \delta[k_x]\delta[k_y]$.
    \item The DFT of the vertical-line signal $f[n_x, n_y] := \delta[n_x]$
    is $F[k_x, k_y] = \frac{1}{N_x} \delta[k_y]$ (a horizontal line).
    \begin{center}
        \frame{\includegraphics[width=0.4\textwidth]{lecture-17/vertical-line-DFT}}
    \end{center}
    \item The DFT of the horizontal-line signal $f[n_x, n_y] := \delta[n_y]$
    is $F[k_x, k_y] = \frac{1}{N_y} \delta[k_x]$ (a vertical line).
    \begin{center}
        \frame{\includegraphics[width=0.4\textwidth]{lecture-17/horizontal-line-DFT}}
    \end{center}
    \item The DFT of an impulse train
    is another, perpendicular impulse train.
    \begin{center}
        \frame{\includegraphics[width=0.4\textwidth]{lecture-17/impulse-train-DFT}}
    \end{center}
\end{itemize}

Here's one more trick:
translation affects DFTs via phase shift.
(Again, recall ``Time Delay'' from
\href{7.html#lecture-and-recitation-properties-of-the-dt-fourier-transform}{Lecture 7}.)

\begin{theorem}{2D DFT Translation}
    Given that $\mathrm{DFT}(f[n_x, n_y]) = F[k_x, k_y]$, we have:
    \[ f'[n_x, n_y] := f[n_x - n_{x, 0}, n_y - n_{y, 0}]
    ~ \implies ~
    \mathrm{DFT}(f'[n_x, n_y]) =
    e^{-jk_x\Omega_{0, x}n_{x, 0}} \cdot e^{-jk_y\Omega_{0, y}n_{y, 0}}
    F[k_x, k_y]. \]
\end{theorem}

See below for the DFT of a translation of a unit-sample signal.
Only the phase changes.

\begin{center}
    \frame{\includegraphics[width=0.4\textwidth]{lecture-17/translation}}
\end{center}

\subsection{Lecture: 2D DFT Visual Properties}

The key idea to keep in mind is the following:
\begin{quote}
    \textbf{Idea.} Lines in the space domain
    manifest as lines perpendicular to their counterparts in the frequency domain.
\end{quote}
Here's an example of the DFT of a triangle.
Notice how the DFT has three linear features,
each perpendicular to a side of the triangle in the original signal.

\begin{center}
    \frame{\includegraphics[width=0.4\textwidth]{lecture-17/triangle-DFT}}
\end{center}

This photo of trees has strong vertical-line features,
so its DFT has a strong horizontal line.

\begin{center}
    \frame{\includegraphics[width=0.4\textwidth]{lecture-17/trees-DFT}}
\end{center}

These are the DFTs of some digits.

\begin{center}
    \frame{\includegraphics[width=0.4\textwidth]{lecture-17/digit-DFT}}
\end{center}

Here's an interesting one:
the DFT of a fingerprint gives this
bright ring of radius $\approx 40$
around a bright dot at the center.

\begin{center}
    \frame{\includegraphics[width=0.4\textwidth]{lecture-17/fingerprint-DFT}}
\end{center}

The reason why is that
a fingerprint locally is just
a collection of evenly-spaced parallel lines.
And a collection of evenly-spaced parallel lines
is exactly what a 2D cosine wave looks like.

But the DFT of a 2D cosine wave is two dots,
antipodal from one another on a circle of radius
equal to the frequency of the cosine wave.
Since a fingerprint consists of evenly-spaced parallel lines
\emph{in all directions}, those pairs of dots
get rotated around this circle in all directions,
forming the bright ring of radius $\approx 40$.

In fact, a common preprocessing step in fingerprint analysis
is to bandpass-filter around the ridge frequency---often implemented
by masking that ring in the DFT before taking the inverse transform,
yielding a high-contrast image of the ridges.

\end{document}